\chapterbackmatter{Solutions to Supplementary Exercises}

\begin{enumerate}
  \SupplementarySolution{1}{Graded Algebra and Supersymmetry Transformations}
  \begin{enumerate}
    \item \emph{The Supertrace Test.}
    For homogeneous block matrices the supertrace has graded cyclicity,
      \[
        \operatorname{str}(XY)
          =(-1)^{|X||Y|}\operatorname{str}(YX).
      \]
      If either matrix is even, this is ordinary cyclicity separately in the
      diagonal blocks.  If both are odd, write
      \(X=\begin{psmallmatrix}0&B\\C&0\end{psmallmatrix}\) and
      \(Y=\begin{psmallmatrix}0&E\\F&0\end{psmallmatrix}\).  Then
      \[
        \operatorname{str}(XY)=\operatorname{tr}(BF)-\operatorname{tr}(CE)
          =-\operatorname{str}(YX),
      \]
      using ordinary cyclicity.  Consequently
      \[
        \operatorname{str}[X,Y\}
          =\operatorname{str}(XY)
           -(-1)^{|X||Y|}\operatorname{str}(YX)=0.
      \]
      The bracket also has degree \(|X|+|Y|\); hence supertraceless matrices
      form a graded subalgebra.
    \item \emph{The Signed Jacobi Identity.}
    Substitute the definition of the bracket twice in each cyclic term.  The
      first term, for example, is
      \[
      \begin{aligned}
        (-1)^{|X||Z|}[X,[Y,Z\}\}
        ={}&(-1)^{|X||Z|}XYZ
           -(-1)^{|X||Z|+|Y||Z|}XZY\\
          &-(-1)^{|X||Y|}YZX
           +(-1)^{|X||Y|+|Y||Z|}ZYX .
      \end{aligned}
      \]
      Expanding the other two cyclic terms produces the same six ordered
      products, each twice with opposite coefficients.  Their sum is zero.
      When all three degrees are one, the inner bracket of each term is an
      anticommutator, while an odd element bracketed with that even result uses
      a commutator.  The identity then reads
      \([X,\{Y,Z\}]+[Y,\{Z,X\}]+[Z,\{X,Y\}]=0\).
    \item \emph{Why a SUSY Variation Is Even.}
    Degrees add modulo two, so
      \(|\epsilon^\alpha Q_\alpha|=1+1=0\), and the adjoint term is even as
      well.  Thus \(G\) is even.  The ordinary commutator with an even operator
      has the same degree as its argument:
      \[
        |[G,X]|=|G|+|X|=|X|.
      \]
      Hence a supersymmetry variation maps bosonic operators to bosonic
      variations and fermionic operators to fermionic variations, although
      \(Q\) itself reverses statistics on states.

      If \(\epsilon\) were commuting, \(\epsilon Q\) would remain odd.
      Moving parameters through odd operators would then miss the Grassmann
      minus signs.  In particular, the commutator of two transformations would
      be represented by the wrong combination of \(QQ^\dagger\) and
      \(Q^\dagger Q\), giving a commutator where the supersymmetry algebra
      requires an anticommutator and therefore failing to close on a
      translation.
  \end{enumerate}

  \SupplementarySolution{2}{Spinorial Lorentz Generators}
The Clifford identity
  \[
    s^\mu\bar s^\nu+s^\nu\bar s^\mu
      =-2\eta^{\mu\nu}\mathbf{1}
  \]
  gives
  \[
  \begin{aligned}
    \Sigma^{\mu\nu}
      &=\frac{i}{4}
        \left(-2\eta^{\mu\nu}\mathbf{1}
              -2s^\nu\bar s^\mu\right)\\
      &=-\frac{i}{2}
        \left(\eta^{\mu\nu}\mathbf{1}
              +s^\nu\bar s^\mu\right).
  \end{aligned}
  \]
  This form is convenient because scalar terms drop out of a commutator.
  Direct multiplication of the Pauli matrices, or one further use of their
  Clifford identity, gives the mostly-plus triple-product formula
  \[
    s^\mu\bar s^\nu s^\rho
      =-\eta^{\mu\nu}s^\rho-\eta^{\nu\rho}s^\mu
       +\eta^{\mu\rho}s^\nu
       -i\epsilon^{\mu\nu\rho\kappa}s_\kappa .
  \]
  Insert this formula twice in
  \[
    [\Sigma^{\mu\nu},\Sigma^{\rho\sigma}]
      =-\frac14
       [s^\nu\bar s^\mu,s^\sigma\bar s^\rho].
  \]
  The epsilon terms cancel between the two matrix orders, whereas the metric
  terms give
  \[
  \begin{aligned}
    [\Sigma^{\mu\nu},\Sigma^{\rho\sigma}]
      =-i\big(&
        \eta^{\nu\rho}\Sigma^{\mu\sigma}
        -\eta^{\nu\sigma}\Sigma^{\mu\rho}\\
        &+\eta^{\mu\sigma}\Sigma^{\nu\rho}
        -\eta^{\mu\rho}\Sigma^{\nu\sigma}\big).
  \end{aligned}
  \]
  Thus \(\Sigma^{\mu\nu}\) obeys precisely the Lorentz
  algebra and acts on an undotted spinor as the \((\tfrac12,0)\)
  representation.

  \SupplementarySolution{3}{The Superspin Casimir}
With the orientation and mostly-plus metric of this volume, take
  \[
    W_\mu=\frac12\epsilon_{\mu\nu\rho\sigma}
                 P^\nu M^{\rho\sigma}.
  \]
  Antisymmetry of \(\epsilon_{\mu\nu\rho\sigma}\), together with the
  Poincar\'e algebra, gives \([W_\mu,P_\rho]=0\).

  For \(x'^\mu=\Lambda^\mu{}_\nu x^\nu+a^\mu\), a left-handed spinor
  transforms as
  \[
    \psi'_a(x')
      =S_{\mathrm L}(\Lambda)_a{}^b\psi_b(x),
    \qquad
    S_{\mathrm L}(\Lambda)
      =\exp\left[-\frac{i}{2}
                   \omega_{\mu\nu}\Sigma^{\mu\nu}\right].
  \]
  The tensor product is
  \[
    \left(\frac12,0\right)\otimes\left(\frac12,0\right)
       =(1,0)\oplus(0,0).
  \]
  Explicitly,
  \[
    \psi_a\chi_b
      =\psi_{(a}\chi_{b)}+\psi_{[a}\chi_{b]}
      =T_{ab}+\epsilon_{ab}S .
  \]
  The symmetric tensor \(T_{ab}\) has three components and furnishes
  \((1,0)\); the antisymmetric tensor has one component and furnishes
  \((0,0)\).  Equivalently, \(S\) is fixed by contracting the last
  equation with \(\epsilon^{ab}\).

  It is useful to denote by \(\mathsf P_{a\dot b}\) the Hermitian momentum
  bispinor.  This notation keeps the signature translation out of the
  dotted-index bookkeeping.  The complete \(N=1\) relations involving
  supercharges are
  \[
  \begin{aligned}
    [Q_a,M^{\mu\nu}]
      &=(\Sigma^{\mu\nu})^b{}_aQ_b,&
    [\overline Q_{\dot a},M^{\mu\nu}]
      &=(\overline\Sigma^{\mu\nu})_{\dot a}{}^{\dot b}
        \overline Q_{\dot b},\\
    [P_\mu,Q_a]&=0,&
    [P_\mu,\overline Q_{\dot a}]&=0,\\
    \{Q_a,Q_b\}&=0,&
    \{\overline Q_{\dot a},\overline Q_{\dot b}\}&=0,\\
    \{Q_a,\overline Q_{\dot b}\}
      &=2\mathsf P_{a\dot b}.&&
  \end{aligned}
  \]
  Here
  \[
    \overline\Sigma^{\mu\nu}
      =\frac{i}{4}
         \left(\bar s^\mu s^\nu-\bar s^\nu s^\mu\right).
  \]
  In the source convention
  \(\mathsf P_{a\dot b}=(s^\mu)_{a\dot b}P_\mu\); the definition above is
  its Hermitian, mostly-plus translation.

  \begin{enumerate}
    \item Since \([W_\mu,P_\rho]=0\) and each supercharge commutes with
      momentum,
      \[
        [B_\mu,P_\rho]
          =[W_\mu,P_\rho]
           -\frac14
             [\overline Q\,\bar s_\mu Q,P_\rho]
          =0.
      \]
    \item It follows at once that
      \[
      \begin{aligned}
        [C_{\mu\nu},P_\rho]
          &=[B_\mu P_\nu-B_\nu P_\mu,P_\rho]\\
          &=[B_\mu,P_\rho]P_\nu
            -[B_\nu,P_\rho]P_\mu=0.
      \end{aligned}
      \]
    \item Put
      \(F_\mu=\overline Q_{\dot b}
        (\bar s_\mu)^{\dot b c}Q_c\).
      Moving \(Q_a\) through the two odd factors gives
      \[
      \begin{aligned}
        [F_\mu,Q_a]
          &=-\{Q_a,\overline Q_{\dot b}\}
              (\bar s_\mu)^{\dot b c}Q_c\\
          &=-2P_\mu Q_a
            -4iP^\nu(\Sigma_{\mu\nu})^b{}_aQ_b .
      \end{aligned}
      \]
      The second line is the vector form of the first, obtained from the
      two-component Clifford identity.  It also displays both the piece
      proportional to \(P_\mu Q_a\) and the spin piece that must cancel
      against the Pauli--Lubanski commutator.
    \item The self-duality relation supplied in the question gives
      \[
        [W_\mu,Q_a]
          =-iP^\nu(\Sigma_{\mu\nu})^b{}_aQ_b .
      \]
      Consequently,
      \[
      \begin{aligned}
        [B_\mu,Q_a]
          &=[W_\mu,Q_a]-\frac14[F_\mu,Q_a]\\
          &=\frac12P_\mu Q_a .
      \end{aligned}
      \]
    \item Using \([P_\mu,Q_a]=0\),
      \[
      \begin{aligned}
        [C_{\mu\nu},Q_a]
          &=[B_\mu,Q_a]P_\nu-[B_\nu,Q_a]P_\mu\\
          &=\frac12
             \left(P_\mu Q_aP_\nu-P_\nu Q_aP_\mu\right)=0.
      \end{aligned}
      \]
      Hermitian conjugation gives
      \([C_{\mu\nu},\overline Q_{\dot a}]=0\).
    \item The operator \(B_\mu\) is a Lorentz vector, so
      \(C_{\mu\nu}\) is an antisymmetric Lorentz tensor.  In particular,
      \[
      \begin{aligned}
        [M^{\rho\sigma},C_{\mu\nu}]
          =i\big(&\delta_\mu^\sigma C^\rho{}_\nu
                 -\delta_\mu^\rho C^\sigma{}_\nu\\
                &+\delta_\nu^\sigma C_\mu{}^\rho
                 -\delta_\nu^\rho C_\mu{}^\sigma\big).
      \end{aligned}
      \]
      Contracting both tensor indices therefore gives
      \[
        [M^{\rho\sigma},\widetilde C_2]
          =[M^{\rho\sigma},C_{\mu\nu}C^{\mu\nu}]=0.
      \]
      Parts (ii) and (v), together with their Hermitian conjugates, show
      that \(\widetilde C_2\) also commutes with \(P_\rho,Q_a\), and
      \(\overline Q_{\dot a}\).  Hence it commutes with every generator
      of the super-Poincar\'e algebra and is a Casimir.
  \end{enumerate}

  \SupplementarySolution{4}{Antisymmetric Central Charges and Their Normal Form}
  \begin{enumerate}
    \item \emph{Antisymmetry of \(Z_{rs}\).}
    Symmetry of the anticommutator under exchange of its two entries yields
      \[
        \epsilon_{\alpha\beta}Z_{rs}
          =\epsilon_{\beta\alpha}Z_{sr}
          =-\epsilon_{\alpha\beta}Z_{sr}.
      \]
      Since \(\epsilon_{\alpha\beta}\) is nonzero for some components,
      \(Z_{rs}=-Z_{sr}\).  For \(N=1\), this says \(Z_{11}=-Z_{11}\), hence
      \(Z_{11}=0\).  At least two supercharges are needed for a scalar central
      charge in this four-dimensional algebra.
    \item \emph{Canonical Blocks of the Central Charge.}
    In the unitary skew-normal basis,
      \[
        Z=\bigoplus_{i=1}^{\lfloor N/2\rfloor}
           z_i\begin{pmatrix}0&1\\-1&0\end{pmatrix}
           \ \oplus\ 0 .
      \]
      Each nonzero block satisfies
      \(Z_i^\dagger Z_i=|z_i|^2\mathbf1_2\).  Thus every
      \(|z_i|^2\) is an eigenvalue of \(Z^\dagger Z\) with multiplicity two;
      for odd \(N\) there is one additional zero eigenvalue.  Consequently
      \[
        \operatorname{Tr}\sqrt{Z^\dagger Z}
          =2\sum_{i=1}^{\lfloor N/2\rfloor}|z_i|.
      \]
  \end{enumerate}

  \SupplementarySolution{5}{The \(N=1\) Algebra, Parity, and State Pairing}
  \begin{enumerate}
    \item Under the usual hypotheses of the theorem, every continuous
      bosonic symmetry of a nontrivial relativistic \(S\)-matrix is a
      direct product of the Poincar\'e group with an internal-symmetry
      group.  Supersymmetry evades the conclusion by weakening the
      assumption that the symmetry generators form an ordinary Lie
      algebra: they instead form a \(\mathbb Z_2\)-graded Lie algebra,
      whose odd--odd bracket is an anticommutator and whose brackets obey
      the graded Jacobi identity.

    \item The Coleman--Mandula conclusion makes an internal generator a
      Lorentz scalar.  The Haag--{\L}opusza\'nski--Sohnius extension then
      permits it to rotate supercharges only through an \(R\)-symmetry.
      Since those symmetries are excluded here,
      \[
        [T_A,Q_a]=[T_A,\overline Q_{\dot a}]=0.
      \]
      The complete simple algebra is
      \[
      \begin{aligned}
        [Q_a,M^{\mu\nu}]
          &=(\Sigma^{\mu\nu})^b{}_aQ_b,\\
        [\overline Q_{\dot a},M^{\mu\nu}]
          &=(\overline\Sigma^{\mu\nu})_{\dot a}{}^{\dot b}
            \overline Q_{\dot b},\\
        [P_\mu,Q_a]&=[P_\mu,\overline Q_{\dot a}]=0,\\
        \{Q_a,Q_b\}
          &=\{\overline Q_{\dot a},\overline Q_{\dot b}\}=0,\\
        \{Q_a,\overline Q_{\dot b}\}
          &=2\mathsf P_{a\dot b}.
      \end{aligned}
      \]
      Here \(\mathsf P_{a\dot b}\) is the Hermitian momentum bispinor;
      it is the mostly-plus translation of
      \((s^\mu)_{a\dot b}P_\mu\) in the source convention.  Also
      \[
        \overline\Sigma^{\mu\nu}
          =\frac{i}{4}
             (\bar s^\mu s^\nu-\bar s^\nu s^\mu).
      \]

      These forms can be recovered rather than guessed.  Lorentz
      covariance fixes the first two lines because \(Q\) and
      \(\overline Q\) furnish the two Weyl representations.  The product
      \((\tfrac12,0)\otimes(0,\tfrac12)\) is a vector, and the only
      allowed vector generator is \(P_\mu\); Hermiticity fixes the
      conjugate relation and a rescaling of \(Q\) fixes the coefficient
      to \(2\).  A scalar term in \(\{Q_a,Q_b\}\) would be a central
      charge antisymmetric in the supersymmetry labels, so it vanishes
      for \(N=1\).  A Lorentz-generator term is incompatible with the
      Jacobi identity involving \(P_\rho\).  The same Jacobi identities
      force the supercharges to commute with momentum.

    \item Write \(P'^\mu=(P^0,-\boldsymbol P)\).  Applying parity to the
      left side of the mixed anticommutator gives
      \[
      \begin{aligned}
        &\widehat{\mathsf P}
          \{Q_a,\overline Q_{\dot b}\}
          \widehat{\mathsf P}^{-1}\\
        &\qquad
          =\eta_{\mathsf P}\eta_{\mathsf P}^{*}
            \,(s^0)_{a\dot c}
            \{\overline Q^{\dot c},Q^d\}
            \,(\bar s^0)_{\dot b d}.
      \end{aligned}
      \]
      The explicit minus sign in the transformation of
      \(\overline Q^{\dot a}\) is cancelled by raising and lowering the
      antisymmetric spinor indices.  Since
      \(|\eta_{\mathsf P}|^2=1\), the phase drops out.  Finally the
      intertwining identity
      \[
        s^0\bar s^\mu s^0
          =\bigl(s^0,-\boldsymbol s\bigr)^\mu
      \]
      sends the momentum bispinor \(\mathsf P\) to the bispinor
      \(\mathsf P'\) constructed from \(P'^\mu\).  Thus
      \[
        \widehat{\mathsf P}
          \{Q_a,\overline Q_{\dot b}\}
          \widehat{\mathsf P}^{-1}
          =2\mathsf P'_{a\dot b},
      \]
      exactly the parity transform of the right side.  The algebra is
      parity consistent.

    \item Fermion parity is defined by
      \[
        (-1)^F\ket{B}=\ket{B},
        \qquad
        (-1)^F\ket{F}=-\ket{F}.
      \]
      Every supercharge reverses statistics, and hence
      \[
        \{(-1)^F,Q_a\}
          =\{(-1)^F,\overline Q_{\dot a}\}=0.
      \]
      Take the trace over a finite-dimensional supermultiplet
      \(\mathcal R_p\) of fixed nonzero four-momentum \(p\).  Cyclicity
      and the preceding anticommutation give
      \[
      \begin{aligned}
        \operatorname{Tr}_{\mathcal R_p}
          \left((-1)^F\{Q_a,\overline Q_{\dot b}\}\right)
          &=
          \operatorname{Tr}_{\mathcal R_p}
          \left(-Q_a(-1)^F\overline Q_{\dot b}
                 +Q_a(-1)^F\overline Q_{\dot b}\right)\\
          &=0.
      \end{aligned}
      \]
      The algebra gives the same trace as
      \[
        2\mathsf p_{a\dot b}
          \operatorname{Tr}_{\mathcal R_p}(-1)^F.
      \]
      For a positive-energy particle multiplet some component of
      \(\mathsf p_{a\dot b}\) is nonzero, and therefore
      \[
        0=\operatorname{Tr}_{\mathcal R_p}(-1)^F
          =n_B-n_F.
      \]
      Zero-energy states annihilated by all supercharges are the familiar
      possible singlet exception to this trace argument.

      If supersymmetry-breaking terms are added explicitly to the
      Lagrangian, the supercharges are no longer conserved:
      \([H,Q_a]\ne0\).  They therefore do not act within a common
      energy eigenspace, and the states are not an exact representation
      of the displayed superalgebra.  The trace step that enforced
      \(n_B=n_F\) is then unavailable.  Calling the breaking terms
      \emph{soft} controls their ultraviolet behavior; it does not
      restore exact state pairing.
  \end{enumerate}

  \SupplementarySolution{6}{Massive \(N=2\) Multiplets and the BPS Bound}
  \begin{enumerate}
    \item The nonzero anticommutators are
      \[
        \{\WeylQ_{ar},\WeylQ_{bs}^\dagger\}
          =2\delta_{rs}\sigma^\mu_{ab}P_\mu,
        \qquad
        \{\WeylQ_{ar},\WeylQ_{bs}\}=e_{ab}Z_{rs},
        \qquad Z_{rs}=-Z_{sr},
      \]
      together with the Hermitian conjugate of the second relation.  The
      charges transform as Weyl spinors,
      \[
        [J_i,\WeylQ_{ar}]
          =-\frac12(\sigma_i)_{ab}\WeylQ_{br},
        \qquad
        [K_i,\WeylQ_{ar}]
          =-\frac{i}{2}(\sigma_i)_{ab}\WeylQ_{br},
      \]
      and
      \[
        [P_\mu,\WeylQ_{ar}]=[P_\mu,\WeylQ_{ar}^\dagger]=0,
        \qquad [Z_{rs},P_\mu]=[Z_{rs},\WeylQ_{at}]=0,
      \]
      with the remaining central-charge commutators and their adjoints also
      zero.

    \item At rest and with \(Z=0\), set
      \[
        a_{ar}=\frac{\WeylQ_{ar}}{\sqrt{2M}},
        \qquad
        a_{ar}^\dagger=\frac{\WeylQ_{ar}^\dagger}{\sqrt{2M}}.
      \]
      These are four canonical fermionic oscillators.  If
      \(a_{ar}\ket\Omega=0\) and \(\ket\Omega\) has spin zero, the states form
      \(\bigwedge^\bullet(\mathbf2_{\rm spin}\otimes\mathbf2_R)\).  By level,
      \[
      \begin{array}{c|c|c}
        F&\text{rotation content}&\text{number of states}\\ \hline
        0&\mathbf1&1\\
        1&2\,\mathbf2&4\\
        2&\mathbf3\oplus3\,\mathbf1&6\\
        3&2\,\mathbf2&4\\
        4&\mathbf1&1
      \end{array}
      \]
      because
      \(\bigwedge^2(\mathbf2\otimes\mathbf2)
       =(\operatorname{Sym}^2\mathbf2\otimes\bigwedge^2\mathbf2)
       \oplus(\bigwedge^2\mathbf2\otimes\operatorname{Sym}^2\mathbf2)\).
      Thus the multiplet contains one spin-one multiplet, four spin-half
      multiplets, and five spin-zero multiplets, for
      \(3+4(2)+5=16=2^4\) states.

    \item Choose the phase so that \(z=|z|\).  With the spinor indices
      identified in the rest frame, define unnormalized combinations
      \[
        A_a=\frac1{\sqrt2}
          (\WeylQ_{a1}+e_{ab}\WeylQ_{b2}^\dagger),
        \qquad
        B_a=\frac1{\sqrt2}
          (\WeylQ_{a1}-e_{ab}\WeylQ_{b2}^\dagger).
      \]
      Direct substitution in the algebra gives
      \[
        \{A_a,A_b^\dagger\}=(2M+|z|)\delta_{ab},
        \qquad
        \{B_a,B_b^\dagger\}=(2M-|z|)\delta_{ab},
        \qquad \{A_a,B_b^\dagger\}=0.
      \]
      Positivity therefore requires
      \[
        M\geq\frac{|Z_{12}|}{2}.
      \]
      At equality, for every \(\ket\psi\),
      \[
        \lVert B_a\ket\psi\rVert^2
          +\lVert B_a^\dagger\ket\psi\rVert^2=0,
      \]
      so both \(B_a\) and \(B_a^\dagger\) act trivially.  Only the two
      normalized \(A_a^\dagger\) creators remain.  On a spin-zero Clifford
      vacuum they give
      \(\mathbf1\oplus\mathbf2\oplus\mathbf1\): two scalar states and one
      spin-half particle, four states in all.  Two of the four real
      supersymmetry pairs annihilate the representation, so it is
      half-BPS.  Relative to the long multiplet, its dimension is reduced by
      a factor of four.

    \item A unitary congruence puts an even-dimensional complex
      antisymmetric matrix into the form
      \[
        UZU^{\mathsf T}
          =\bigoplus_{i=1}^{N/2}
            z_i\begin{pmatrix}0&1\\-1&0\end{pmatrix},
        \qquad z_i\geq0.
      \]
      Each block repeats the preceding construction and yields
      \(2M\pm z_i\) as the two oscillator norms.  Hence
      \[
        M\geq\frac12\max_i z_i.
      \]
      A generic massive multiplet has \(2N\) creation operators and therefore
      \(2^{2N}\) states for a one-dimensional Clifford vacuum.  Each
      saturated skew block makes two creation operators null.  If exactly
      \(k\) blocks saturate, the remaining exterior algebra has dimension
      \[
        2^{2N-2k}.
      \]
      For a Clifford vacuum of spin \(j\), this is multiplied by
      \(2j+1\).  The shortening records the corresponding preserved
      supercharges; in particular, saturation of every block is half-BPS.
  \end{enumerate}

  \SupplementarySolution{7}{Massive \(N=1\) Multiplets and Parity}
  \begin{enumerate}
    \item \emph{A Long Massive \(N=1\) Multiplet.}
    The zero-creation level is one spin-\(j\) representation.  Because the
      creation operators transform as spin \(1/2\), the one-creation level is
      \[
        j\otimes\frac12
          =\left(j+\frac12\right)\oplus
           \left(j-\frac12\right).
      \]
      Antisymmetry leaves only the spin-singlet product of the two creation
      operators, so the two-creation level is a second spin-\(j\)
      representation.  The full content is therefore
      \[
        \left(j+\frac12\right)\oplus j\oplus j
          \oplus\left(j-\frac12\right).
      \]
      The two even levels contain \(2(2j+1)=4j+2\) spin states.  The odd level
      contains \((2j+2)+(2j)=4j+2\).  Which side is bosonic depends on whether
      the Clifford vacuum has integral or half-integral spin, but the equality
      is the same.
    \item \emph{The Collapsed Massive Multiplet.}
    For \(j=0\), the tensor product
      \(0\otimes1/2\) contains only spin \(1/2\); a negative-spin
      representation has no meaning and is absent.  The zero- and two-creation
      levels are each spin zero, while the one-creation level is spin \(1/2\).
      Thus the spectrum is
      \[
        0\oplus0\oplus\frac12.
      \]
      The two scalars give two physical spin states in total, and the massive
      spinor gives two, so bosonic and fermionic degrees of freedom again
      match.
    \item \emph{Parity Along the Massive Ladder.}
    Parity exchanges annihilation and creation operators.  It therefore
      exchanges the spin-\(j\) zero-creation state \(\ket{j,m}\) with the
      two-creation state
      \(\ket{j,m}^{\mathrm b}\propto
        Q^\dagger_{1/2}Q^\dagger_{-1/2}\ket{j,m}\).
      With the stated convention, applying parity twice through the two odd
      operators gives the relative action
      \[
        P\ket{j,m}=\eta\ket{j,m}^{\mathrm b},\qquad
        P\ket{j,m}^{\mathrm b}=-\eta\ket{j,m},
      \]
      after fixing the common phase by assigning parity \(\eta\) to the
      one-creation multiplets.  Diagonalizing this \(2\times2\) action gives
      \[
        \frac{\ket{j,m}\pm i\ket{j,m}^{\mathrm b}}{\sqrt2},
        \qquad P=\pm i\eta .
      \]
      The Clebsch--Gordan combinations at the one-creation level, of spins
      \(j\pm1/2\), both retain parity \(\eta\).  At \(j=0\), the two scalars
      have parities \(\pm i\eta\), while the spin-\(1/2\) particle has parity
      \(\eta\).
  \end{enumerate}

  \SupplementarySolution{8}{Massless Oscillators, Helicity, and CPT}
  \begin{enumerate}
    \item \emph{The Massless-Frame Oscillator.}
    In the chosen null frame the anticommutator matrix has eigenvalues
      \(4E\) and \(0\):
      \[
        \{Q_\alpha,Q^\dagger_{\dot\beta}\}
          =4E\begin{pmatrix}1&0\\0&0\end{pmatrix}_{\alpha\dot\beta}
      \]
      after a possible relabeling of components.  For every state
      \(\ket\psi\),
      \[
        \lVert Q_2\ket\psi\rVert^2+
        \lVert Q_2^\dagger\ket\psi\rVert^2=0.
      \]
      Unitarity therefore makes both \(Q_2\) and \(Q_2^\dagger\) act trivially.
      The normalized nonzero pair
      \(a=Q_1/\sqrt{4E}\), \(a^\dagger=Q_1^\dagger/\sqrt{4E}\)
      obeys \(\{a,a^\dagger\}=1\).  In the highest-helicity construction,
      \(a^\dagger\) lowers helicity by \(1/2\), and its adjoint \(a\) raises it
      by \(1/2\).
    \item \emph{The Exterior Algebra of Helicity.}
    The \(N\) creation operators anticommute, so a basis at level \(k\) is
      \[
        a^\dagger_{i_1}\cdots a^\dagger_{i_k}\ket\Omega,
        \qquad i_1<\cdots<i_k.
      \]
      This is \(\bigwedge^k\mathbb C^N\), of dimension \(\binom Nk\).
      Every creation operator lowers helicity by \(1/2\), giving
      \(\lambda_k=\lambda_{\max}-k/2\).  Summing all levels,
      \[
        \dim\mathcal H=\sum_{k=0}^N\binom Nk=(1+1)^N=2^N.
      \]
    \item \emph{The CPT Self-Conjugacy Test.}
    CPT reverses every helicity, so the interval
      \([\lambda_{\max}-N/2,\lambda_{\max}]\) is invariant only if its endpoints
      are negatives:
      \[
        \lambda_{\max}-\frac N2=-\lambda_{\max},
        \qquad\text{or}\qquad
        \lambda_{\max}=\frac N4.
      \]
      The binomial multiplicities are then also symmetric because
      \(\binom Nk=\binom N{N-k}\).  If this condition fails, CPT produces a
      second multiplet whose top helicity is the negative of the original
      bottom helicity,
      \[
        \lambda_{\max}^{\mathrm{CPT}}
          =-\lambda_{\max}+\frac N2.
      \]
  \end{enumerate}

  \SupplementarySolution{9}{The \(N=1\) Chiral and Vector Multiplets}
  \begin{enumerate}
    \item \emph{The Massless Chiral Multiplet.}
    Beginning with helicity \(+1/2\), the single \(N=1\) creation operator
      gives
      \[
        \left(+\frac12,\,0\right).
      \]
      CPT supplies \((0,-1/2)\).  The two helicity-zero states are the particle
      and antiparticle states of one complex scalar, while the
      \(\pm1/2\) states are the two on-shell states described by one Weyl
      fermion and its conjugate.  Thus there are two real bosonic and two
      fermionic on-shell degrees of freedom.
    \item \emph{The Massless Vector Multiplet.}
    The multiplet beginning at \(+1\) contains
      \((+1,+1/2)\).  Its CPT conjugate contains
      \((-1,-1/2)\).  The \(\pm1\) states are the two transverse polarizations
      of a massless vector, and the \(\pm1/2\) states form a Majorana gaugino
      (equivalently, one Weyl gaugino with its conjugate).  Each side has two
      physical polarizations, as supersymmetry requires.
  \end{enumerate}

  \SupplementarySolution{10}{Extended Multiplets and the \(N=4\) \(N=1\) Decomposition}
  \begin{enumerate}
    \item \emph{Field Content in \(N=1\) Language.}
    Relative to a chosen \(N=1\) subalgebra, the \(N=4\) multiplet is
      \[
        V=(A_\mu,\lambda,D),\qquad
        \Phi_i=(\phi_i,\psi_i,F_i),\quad i=1,2,3,
      \]
      with every field in the adjoint representation.  The vector contributes
      one gauge boson and one Weyl fermion; the three chirals contribute three
      complex scalars and three more Weyl fermions.  Thus the decomposition
      contains the expected one vector, four Weyl fermions, and six real
      scalars.
    \item \emph{The \(N=1\) Superspace Action.}
    Up to a common convention for rescaling the adjoint chirals, the action is
      \[
      \begin{aligned}
        S=\frac1{g^2}\int d^4x\,\operatorname{Tr}\Bigg\{&
          \int d^4\theta\,
             \sum_{i=1}^3
             \Phi_i^\dagger e^{2V}\Phi_i e^{-2V}\\
          &+\left[
          \int d^2\theta\left(
             \frac14 W^\alpha W_\alpha
             +\frac{\sqrt2}{3!}\epsilon^{ijk}
                \Phi_i[\Phi_j,\Phi_k]\right)
          +\mathrm{H.c.}\right]\Bigg\}.
      \end{aligned}
      \]
      The coefficient of the cubic superpotential is tied to the gauge
      coupling; choosing it independently would retain only \(N=1\).
      The \(D\)-equation produces
      \(\sum_i[\phi_i,\phi_i^\dagger]\), while
      \(F_i^\dagger\) is proportional to
      \(\epsilon_{ijk}[\phi_j,\phi_k]\).  Writing
      \(\phi_i=(X^{2i-1}+iX^{2i})/\sqrt2\) therefore combines both
      contributions into
      \[
        V_{\rm scalar}
          =\frac1{4g^2}\sum_{I,J=1}^6
             \operatorname{Tr}
             \big([X^I,X^J]^\dagger[X^I,X^J]\big),
      \]
      the six-scalar commutator potential.
    \item \emph{The Hidden \(SO(6)\) Symmetry.}
    The six \(X^I\) transform as the vector \(\mathbf6\) of
      \(SO(6)\), while
      \((\lambda,\psi_1,\psi_2,\psi_3)\) form the
      \(\mathbf4\) of
      \(SU(4)\simeq\operatorname{Spin}(6)\).  The Yukawa matrices are the
      chiral gamma matrices of \(SO(6)\), so the Yukawa couplings and the
      commutator potential are invariant under the full group.  The
      \(N=1\) presentation makes only the subgroup preserving the selected
      supercharge manifest.

      If the full \(SU(4)_R\) is an exact symmetry, conjugating the manifest
      supercharge by it generates a \(\mathbf4\) of supercharges.  Their
      conjugates and algebra therefore furnish all sixteen real
      supercharges.  Hence exact \(N=1\) supersymmetry plus this full
      \(R\)-symmetry upgrades the theory to \(N=4\); merely retaining the
      manifest \(SU(3)\times U(1)\) subgroup would not.
  \end{enumerate}

  \SupplementarySolution{11}{The \(N=8\) Gravity Multiplet and Helicity Cutoffs}
  \begin{enumerate}
    \item \emph{The \(N=8\) Gravity Multiplet.}
    The binomial construction gives
      \[
      \begin{array}{c|rrrrrrrrr}
        \lambda&2&\frac32&1&\frac12&0&
          -\frac12&-1&-\frac32&-2\\ \hline
        \text{multiplicity}&1&8&28&56&70&56&28&8&1 .
      \end{array}
      \]
      Hence the multiplet contains one graviton, eight gravitini, 28 vectors,
      56 spin-\(1/2\) particles, and 70 real scalar states.  The table is
      invariant under \(\lambda\mapsto-\lambda\), so it is CPT
      self-conjugate.
    \item \emph{Helicity Cutoffs on Extended SUSY.}
    A CPT-self-conjugate multiplet has
      \(\lambda_{\max}=N/4\) and
      \(\lambda_{\min}=-N/4\).  If all helicities satisfy
      \(|\lambda|\leq1\), then \(N/4\leq1\), hence \(N\leq4\).  If gravity is
      allowed but \(|\lambda|\leq2\), the same argument gives \(N\leq8\).
      The first bound concerns perturbative interacting global supersymmetry
      with no massless higher-spin particles.  The second concerns local
      supersymmetry with a graviton but no massless fields above spin two.
      Dropping a particle interpretation or allowing free higher-spin systems
      evades the corresponding physical premise.
  \end{enumerate}

  \SupplementarySolution{12}{Protected Helicity Sums and Central-Charge Parity}
  \begin{enumerate}
    \item \emph{Helicity Supertraces.}
    Expanding near \(t=0\),
      \[
        e^{t\lambda_{\max}}(1-e^{-t/2})^N
          =e^{t\lambda_{\max}}
           \left(\frac t2+O(t^2)\right)^N
          =\frac{t^N}{2^N}+O(t^{N+1}).
      \]
      But the \(m\)th derivative at the origin is exactly
      \(B_m=\sum(-1)^F\lambda^m\).  All derivatives of order \(m<N\) vanish,
      while
      \[
        B_N=\frac{N!}{2^N}.
      \]
      This sign assumes the highest-helicity state is bosonic, as in the given
      generating function; a fermionic Clifford vacuum multiplies every
      supertrace by \(-1\).
    \item \emph{Parity and the Central-Charge Phase.}
    Transforming both supercharges gives
      \[
      \begin{aligned}
        P^{-1}\{Q_{\alpha r},Q_{\beta s}\}P
          &=-\epsilon_{\alpha\gamma}\epsilon_{\beta\delta}
            \{Q^\dagger_{\dot\gamma r},
              Q^\dagger_{\dot\delta s}\}\\
          &=-\epsilon_{\alpha\beta}Z^\dagger_{rs}.
      \end{aligned}
      \]
      Comparison with the original algebra therefore yields
      \[
        P^{-1}Z_{rs}P=-Z^\dagger_{rs}.
      \]
      If \(Z_{rs}\) has eigenvalue \(z_{rs}\), parity maps that sector to the
      one with eigenvalue \(-z_{rs}^*\).  A single sector can carry parity
      eigenstates only when it is invariant under this map (for one charge
      block, \(z=-z^*\) in the stated convention); otherwise the parity
      representation must include and exchange the two sectors.  Rephasing
      \(Q_r\) changes the phase of \(Z_{rs}\) twice as fast and simultaneously
      changes the phase in the parity law.  Thus ``purely imaginary'' is
      convention-dependent, while the requirement that parity preserve or
      pair central-charge sectors is invariant.
  \end{enumerate}

\end{enumerate}
